\chapter{Циклически-изотипический проектор полного меню рёбер}
\label{ch:v7-rooted-cycle-isotypic-edge-projector-gate}

\section{Два независимо выведенных признака}

Предыдущий гейт выделил четыре новых ребра единственного примитивного цикла,
укоренённого в исходном $H_{15}$. Две требуемые вектороподобные массы в этот
цикл не входят. Проверим, дополняет ли циклический признак независимый
представленческий признак без использования меток ``старое'' и ``новое''.

Пусть $\mathcal E_{\rm new}$ --- прямая сумма всех одиннадцати новых
первопорядково допустимых пространств стрелок:
\begin{equation}
 \mathcal E_{\rm new}
 =\bigoplus_{e\in E_{\rm new}}
 \operatorname{Hom}(\mathcal H_{s(e)},\mathcal H_{t(e)}),
 \qquad |E_{\rm new}|=11.
 \label{eq:v7-edge-projector-space}
\end{equation}
Каждое слагаемое является полным калибровочно-ковариантным блоком, а не
отдельной матричной компонентой.

\section{Проектор циклического участия}

Для единственного укоренённого монома $\mathcal R_6(x)$ определим формальную
степень ребра
\begin{equation}
 \nu_e=x_e\frac{\partial}{\partial x_e}\log\mathcal R_6(x)
 \in\{0,1\}.
 \label{eq:v7-edge-projector-cycle-degree}
\end{equation}
Поскольку $\mathcal R_6$ является одним квадратсвободным мономом, $\nu_e$
не зависит от значения ненулевых переменных и просто измеряет участие ребра
в цикле. На новом пространстве стрелок получаем ортогональный проектор
\begin{equation}
 P_C=\sum_{e\in E_{\rm new}}\nu_e|e\rangle\langle e|,
 \qquad \rank P_C=4.
 \label{eq:v7-edge-projector-cycle}
\end{equation}
Его опора равна
\begin{equation}
 E_C=\{L_LY_R,Q_LY_R,X_Le_R,X_Lu_R\}.
 \label{eq:v7-edge-projector-cycle-support}
\end{equation}
Ложноположительных рёбер нет, но отсутствуют $X_LX_R$ и $Y_LY_R$.

\section{Изотипический проектор}

Ребро назовём изотипическим, если его левый и правый концы несут одно и то
же представление $SU(3)\times SU(2)\times U(1)_Y$ после забывания
хиральности:
\begin{equation}
 \epsilon_e^{\rm iso}
 =\mathbf1\!\left[\rho_{s(e)}\simeq\rho_{t(e)}\right].
 \label{eq:v7-edge-projector-isotypic-indicator}
\end{equation}
Такое условие является условием калибровочно-инвариантного
вектороподобного массового блока. Диаграммы Краевского именно связывают
вершины с хиральными мультиплетами, а стрелки с допустимыми массовыми и
юкавскими блоками \cite{v7-edge-projector-krajewski,v7-edge-projector-jureit}.

Соответствующий проектор имеет вид
\begin{equation}
 P_I=\sum_{e\in E_{\rm new}}\epsilon_e^{\rm iso}|e\rangle\langle e|,
 \qquad \rank P_I=4,
 \label{eq:v7-edge-projector-isotypic}
\end{equation}
и опору
\begin{equation}
 E_I=\{L_LY_R,Y_LY_R,X_Le_R,X_LX_R\}.
 \label{eq:v7-edge-projector-isotypic-support}
\end{equation}
Она включает обе недостающие массы. Ложноположительных рёбер снова нет;
пропущены только два смешанных ребра $Q_LY_R$ и $X_Lu_R$.

\section{Точное объединение шести рёбер}

Оба проектора диагональны в каноническом разложении по пространствам
стрелок, поэтому коммутируют. Их пересечение имеет ранг два. Проектор на
объединение равен
\begin{equation}
 P_*=P_C\vee P_I=P_C+P_I-P_CP_I,
 \qquad P_*^2=P_*,\qquad \rank P_*=6.
 \label{eq:v7-edge-projector-union}
\end{equation}
Полный конечный аудит даёт
\begin{align}
 \operatorname{supp}P_*
 &=\{L_LY_R,Q_LY_R,X_LX_R,X_Le_R,X_Lu_R,Y_LY_R\},
 \nonumber\\
 \operatorname{supp}(I-P_*)
 &=\{L_LX_R,X_LY_R,X_Ld_R,X_RY_L,Y_Le_R\}.
 \label{eq:v7-edge-projector-exact-partition}
\end{align}
Первая строка совпадает ровно с шестью рёбрами целевого ремонта, вторая ---
с пятью ранее нежелательными. Ни один относительный вес между циклическим и
изотипическим признаками не вводится: формула объединения двух коммутирующих
проекторов единственна.

\section{Каноническая знаковая градуировка}

Проектор автоматически задаёт инволюцию пространства новых стрелок
\begin{equation}
 \Gamma_E=I-2P_*,
 \qquad \Gamma_E^*=\Gamma_E,
 \qquad \Gamma_E^2=I.
 \label{eq:v7-edge-projector-grading}
\end{equation}
Её квадратичная форма равна
\begin{equation}
 q_E(z)=\langle z,\Gamma_Ez\rangle
 =\sum_{e\notin E_*}\|z_e\|^2-\sum_{e\in E_*}\|z_e\|^2.
 \label{eq:v7-edge-projector-quadratic-form}
\end{equation}
Тем самым все шесть целевых блоков получают отрицательный знак запуска, а
пять конкурирующих --- положительную щель. Сигнатуры гессиана равны
\begin{equation}
 (n_-,n_0,n_+)_{\rm one\ gen}=(12,0,10),
 \qquad
 (n_-,n_0,n_+)_{3\times3\rm\ family}=(108,0,90).
 \label{eq:v7-edge-projector-hessian-signatures}
\end{equation}
Знак сохраняется на ориентированном Real-сопряжении ребра, поскольку
$P_*$ определён на полном неориентированном блоке.

\section{Оставшийся родительский разрыв}

Полученная $\Gamma_E$ является канонической градуировкой пространства
стрелок относительно уже установленного фона $H_{15}$. Но её существование
ещё не доказывает, что она входит в исходную суперсвязность или спектральное
действие. Квадратичная форма \eqref{eq:v7-edge-projector-quadratic-form}
сама по себе не ограничена снизу. Не выведены:
\begin{equation}
 \begin{gathered}
  \text{общий масштаб запуска},\qquad
  \text{единая квартичная стабилизация},\\
  \text{ненулевой минимум всех шести целевых блоков}.
 \end{gathered}
 \label{eq:v7-edge-projector-parent-gaps}
\end{equation}

\begin{theorem}[Точный проектор шести рёбер]
Объединение проектора участия в единственном $H_{15}$-укоренённом
примитивном цикле и проектора изотипических лево-правых блоков совпадает
ровно с шестью рёбрами целевого четырёхвершинного ремонта. Производная
инволюция $\Gamma_E=I-2P_*$ даёт отрицательный квадратичный знак всем шести
целевым и положительный знак всем пяти нежелательным блокам без
относительных коэффициентов и без меток старой/новой копии.
\end{theorem}

\begin{nogocriterion}[Проектор ещё не равен действию]
Нельзя объявлять знаковую форму полным физическим родителем, пока одна
Real-полевая суперсвязность не выводит одновременно $\Gamma_E$, общий
масштаб и стабилизирующий член. Независимое добавление
$(\|P_*z\|^2-\mu^2)^2$, отдельных масс рёбер или ручных юкавских отношений
запрещено.
\end{nogocriterion}

\begin{protocol}
\begin{enumerate}
 \item размерность пространства новых комплексных рёбер: \textbf{$11$};
 \item ранг циклического проектора: \textbf{$4$};
 \item ранг изотипического проектора: \textbf{$4$};
 \item ранг пересечения: \textbf{$2$};
 \item ранг объединения: \textbf{$6$};
 \item ранг дополнения: \textbf{$5$};
 \item объединение совпадает с целевым меню: \textbf{да};
 \item ложноположительных рёбер: \textbf{$0$};
 \item свободный относительный вес: \textbf{нет};
 \item сигнатура вещественного гессиана на поколение: \textbf{$(12,0,10)$};
 \item семейная сигнатура: \textbf{$(108,0,90)$};
 \item Real-парность знаков: \textbf{сохранена};
 \item общий масштаб и квартичная стабилизация: \textbf{не выведены};
 \item итог: \textbf{точный канонический проектор, родитель действия открыт}.
\end{enumerate}
\end{protocol}

Машинный сертификат сохранён в
\path{s2t_v7_rooted_cycle_isotypic_edge_projector_gate_results.json}.

\paragraph{Уровень подсчёта сигнатур.}
Сигнатуры в \eqref{eq:v7-edge-projector-hessian-signatures} относятся к
пространству меток рёбер: одно комплексное число на ребро для одного
поколения и одна комплексная матрица $3\times3$ на ребро на семейном уровне.
Если раскрыть внутренние калибровочные компоненты полного блока, их
кратности дополнительно умножаются на размерность соответствующего
мультиплета, но знак, заданный $\Gamma_E$, наследуется всеми компонентами и
не меняется.

\begin{thebibliography}{9}
\bibitem{v7-edge-projector-krajewski}
T.~Krajewski,
\emph{Classification of Finite Spectral Triples},
Journal of Geometry and Physics 28 (1998), 1--30;
arXiv:hep-th/9701081.
\bibitem{v7-edge-projector-jureit}
J.-H.~Jureit, T.~Schuecker, C.~A.~Stephan,
\emph{On a Classification of Irreducible Almost-Commutative Geometries V},
Journal of Mathematical Physics 50 (2009), 072303; arXiv:0901.3214.
\end{thebibliography}